% Provenance (never printed): see sources/outline_sources.md (Year 2: radiative forcing
% quantified; feedbacks and climate sensitivity; detection and attribution; scenarios; the
% remaining carbon budget; components of sea-level rise; physical science only). Every statement
% from IPCC AR6 WG1 (SPM; chapters 3, 5, 7 and 9; Annex III data), ledger rows in
% sources/ledger/book2.md (ipcc:*) and sources/ledger/book4.md (ar6*:*). Curves and numbers from
% figdata/bachelor-2/anthropogenic.py and the Book 4 second-writer scratch calculator calc/ch27.py.
\chapter{Anthropogenic Climate Change}\label{ch:b2:anthropogenic-climate-change}

Since 1970 the global surface has warmed faster than in any other fifty-year period of at least
the last two thousand years (\emph{high confidence}). The school volumes told what has changed and
why. This chapter turns the account into a budget: how many watts per square metre the human
drivers add, how many the warmer planet sends back to space, how much warming each tonne of carbon
dioxide buys, and how the rising sea is shared between warming water and melting ice. Each number
comes with its uncertainty, in the calibrated language of the assessments.
% ledger: ipcc:fastest-2000yr, ipcc:likelihood

\begin{recall}
\textbf{From the school volumes}: radiative forcing, $\qty{3.93}{W.m^{-2}}$ for a doubling of
carbon dioxide, and its logarithmic law; climate feedback; \omterm{def:b2:anthropogenic-climate-change:ecs}{equilibrium climate sensitivity}, best
estimate \qty{3}{\celsius}, \emph{likely} 2.5 to \qty{4}{\celsius}; \omterm{def:b2:anthropogenic-climate-change:da}{attribution}, emission scenario
and projection; the calibrated language: \emph{likely} 66 to \qty{100}{\percent}, \emph{very
likely} 90 to \qty{100}{\percent}, a range $[x \text{ to } y]$ being a very likely range. Human
influence has warmed the climate by \qty{1.07}{\celsius} between 1850--1900 and 2010--2019.
\textbf{From the physics books}: the black body and the Stefan--Boltzmann law; thermal expansion.
\textbf{From the biology books}: the carbon cycle, the airborne fraction, carbon sinks and the
carbon budget of the biosphere. \textbf{From \cref{ch:b2:deep-time-climate}}: the slow carbon cycle,
which takes back an excess of carbon dioxide only over hundreds of thousands of years.
\end{recall}
% ledger: ipcc:erf-2xco2, ipcc:log-co2, ipcc:ecs, ipcc:likelihood, ipcc:attribution

\section{Forcing quantified}\label{sec:b2:anthropogenic-climate-change:forcing}

\begin{definition}[Effective radiative forcing]\label{def:b2:anthropogenic-climate-change:erf}
The \emph{effective radiative forcing}\index{effective radiative forcing} (ERF) of a perturbation
--- a change in a greenhouse gas, in aerosols, in the Sun --- is the change it causes in the net
downward flux of radiation at the top of the atmosphere, in \unit{W.m^{-2}}, once the atmosphere,
the clouds and the land surface have adjusted to it, but before the global surface temperature
has responded. The adjustments counted are those caused by the agent itself, independently of the
surface temperature.
\end{definition}
% ledger: ar6ch7:erf-def

The definition draws the line between forcing and response by what depends on the surface
temperature, not by speed. The instantaneous forcing of carbon dioxide is the change of flux just
after the gas is added; the stratosphere, which then cools within months, changes it a first time;
the troposphere, its water vapour and its clouds change it again. For a doubling of carbon
dioxide the assessment finds \qty{3.75}{W.m^{-2}} once the stratosphere has adjusted, and
adjustments in the troposphere that partly cancel: \qty{-0.60}{W.m^{-2}} from its warming, +0.22
from water vapour, +0.45 from clouds and +0.11 from the land surface, \qty{+0.18}{W.m^{-2}} in all.
The ERF is $3.75 + 0.18 = \qty{3.93}{W.m^{-2}}$, \qty{6}{\percent} more than the
\qty{3.71}{W.m^{-2}} of the previous assessment (\cref{exo:b2:anthropogenic-climate-change:1}).
% ledger: ar6ch7:2xco2, ar6ch7:erf-def

The forcing of every agent since 1750 has been reconstructed year by year. The figure below shows
it: carbon dioxide at \qty{2.16}{W.m^{-2}} in 2019, the other greenhouse gases --- methane,
nitrous oxide, the halogenated gases and ozone --- at \qty{1.69}{W.m^{-2}} together, and the
aerosols at \qty{-1.06}{W.m^{-2}}. The total human forcing is \qty{2.72}{W.m^{-2}}, and its very
likely range, 1.96 to \qty{3.48}{W.m^{-2}}, is wide mainly because of the aerosols. Their ERF over
1750--2014 is assessed at $-1.3$ [$-2.0$ to $-0.6$] \unit{W.m^{-2}} (\emph{medium confidence}),
most of it through their effect on clouds, which brighten when aerosols supply more droplets.
Volcanic eruptions add short negative spikes (\cref{ch:b2:eruptive-styles}): the deepest in the
record, \qty{-4.8}{W.m^{-2}} in 1816, follows the eruption of Tambora in 1815.
% ledger: data:b4-ar6-erf, ipcc:erf-total, ipcc:erf-agents, ar6ch7:aerosol, usgs:msh-comparisons

\begin{omfigure}
\centering
\begin{tikzpicture}
\begin{axis}[omrecord, width=0.92\linewidth, height=6.2cm, xmin=1750, xmax=2020, ymin=-5.2,
  ymax=3.8, xlabel={year}, ylabel={ERF since 1750 (\unit{W.m^{-2}})},
  xtick={1750,1800,...,2000}, /pgf/number format/1000 sep={}, xtick pos=bottom, ytick pos=left,
  tick label style={font=\scriptsize}, label style={font=\small},
  legend style={at={(0.5,-0.24)}, anchor=north, font=\scriptsize, draw=black!30,
  legend columns=3, /tikz/every even column/.append style={column sep=6pt}},
  legend cell align=left]
% computed in figdata/bachelor-2/anthropogenic.py part erf; ledger: data:b4-ar6-erf
\addplot[name path=hi, draw=none, forget plot] table[x=year, y=hi] {figdata/out/bachelor-2-anthropogenic-erf.dat};
\addplot[name path=lo, draw=none, forget plot] table[x=year, y=lo] {figdata/out/bachelor-2-anthropogenic-erf.dat};
\addplot[black!15, forget plot] fill between[of=hi and lo];
\addplot[black!45, line width=0.4pt] table[x=year, y=natural] {figdata/out/bachelor-2-anthropogenic-erf.dat};
\addlegendentry{natural (volcanoes, Sun)}
\addplot[omThm, line width=1pt] table[x=year, y=co2] {figdata/out/bachelor-2-anthropogenic-erf.dat};
\addlegendentry{carbon dioxide}
\addplot[omThm!50!orange, line width=1pt, densely dashed] table[x=year, y=other_ghg] {figdata/out/bachelor-2-anthropogenic-erf.dat};
\addlegendentry{other greenhouse gases and ozone}
\addplot[omDef, line width=1pt] table[x=year, y=aerosol] {figdata/out/bachelor-2-anthropogenic-erf.dat};
\addlegendentry{aerosols}
\addplot[black, line width=1.4pt] table[x=year, y=anthro] {figdata/out/bachelor-2-anthropogenic-erf.dat};
\addlegendentry{total human (grey: 5--95\,\%)}
\draw[black!40] (axis cs:1750,0) -- (axis cs:2020,0);
\node[anchor=west, font=\scriptsize, text=black!60, inner sep=1pt] at (axis cs:1817.5,-4.75) {$\leftarrow$ Tambora};
\end{axis}
\end{tikzpicture}
\par\smallskip{\small \omterm{def:b2:anthropogenic-climate-change:erf}{Effective radiative forcing} from 1750 to 2019, relative to 1750, by driver.
The total human forcing reaches \qty{2.72}{W.m^{-2}} in 2019, with the very likely range shaded.
Data: IPCC Sixth Assessment Report, Working Group~I, Annex~III extended data (CC BY 4.0).}
\end{omfigure}
% values: anthro 2.7204 (1.9569-3.4819), co2 2.1563, other_ghg 1.6850, aerosol -1.0583, volcanic -4.8162 in 1816 (anthropogenic.py numbers)

\section{The energy budget}\label{sec:b2:anthropogenic-climate-change:energy}

A positive forcing makes the planet gain energy. As it warms it radiates more to space, and the
gain shrinks; the climate settles only when the two balance.

\begin{definition}[Earth's energy imbalance; climate feedback parameter]\label{def:b2:anthropogenic-climate-change:eei}
\emph{Earth's energy imbalance}\index{Earth's energy imbalance} $N$ is the net flux of energy into
the climate system at the top of the atmosphere, averaged over the globe and over a year or more,
in \unit{W.m^{-2}}: absorbed sunlight minus outgoing infrared. It is the rate at which the ocean,
the land, the ice and the air gain heat, per square metre of the Earth's surface.
\end{definition}

\begin{definition}[Climate feedback parameter]\label{def:b2:anthropogenic-climate-change:lambda}
The \emph{climate feedback parameter}\index{climate feedback parameter} $\lambda$ is the change in
the net downward flux at the top of the atmosphere per degree of global surface warming, in
\unit{W.m^{-2}.\celsius^{-1}}, caused by the response of the climate system to that warming. A
stable climate has $\lambda < 0$: warming increases the loss of energy to space.
\end{definition}

The school volume wrote the balance with the size of $\lambda$, a positive number; from here on the
sign is kept, as in the assessments, which write the parameter $\alpha$.
% ledger: ar6ch7:energy-budget

\begin{proposition}[Global energy balance]\label{prop:b2:anthropogenic-climate-change:balance}
Let $\Delta F$ be the \omterm{def:b2:anthropogenic-climate-change:erf}{effective radiative forcing} and $\Delta T$ the global surface warming
relative to an equilibrium state. If the response of the outgoing flux is linear in $\Delta T$,
\[
  \Delta N = \Delta F + \lambda\,\Delta T .
\]
At a new equilibrium $\Delta N = 0$, so $\Delta T_{\text{eq}} = -\Delta F/\lambda$.
\end{proposition}
\begin{proof}
By definition the forcing changes the net downward flux by $\Delta F$ at fixed surface
temperature, and the response to the warming changes it by $\lambda\,\Delta T$ to first order;
the imbalance is their sum. Setting it to zero gives the equilibrium.
\end{proof}

Measuring $N$ means measuring the heat the system gains, and the ocean holds nearly all of it.
From 1971 to 2018 the climate system gained \qty{435}{ZJ} (zettajoules, \qty{e21}{J}): the ocean
\qty{91}{\percent}, \qty{242}{ZJ} in its top \qty{700}{m} and \qty{154}{ZJ} below; the land
\qty{5}{\percent}, the melting ice \qty{3}{\percent}, the atmosphere \qty{1}{\percent}. The mean
imbalance was \qty{0.57}{W.m^{-2}} over 1971--2018, and it has grown: \qty{0.79}{W.m^{-2}}
[0.52 to 1.06] over 2006--2018 (\cref{exo:b2:anthropogenic-climate-change:2}).
% ledger: ar6ch7:inventory, ipcc:heating

\begin{omfigure}
\centering
\begin{tikzpicture}[font=\small, >=Stealth]
\draw[fill=omDef!8, draw=omDef!60] (-3.2,0) rectangle (3.2,1.6);
\node[align=center] at (0,1.05) {climate system};
\node[align=center, font=\scriptsize] at (0,0.5) {ocean \qty{91}{\percent}, land \qty{5}{\percent}, ice \qty{3}{\percent}, air \qty{1}{\percent}};
\draw[->, line width=2.2pt, omThm] (-2.0,3.3) -- node[pos=0.12, left=3pt, align=right, font=\scriptsize] {forcing $\Delta F$} (-2.0,1.7);
\draw[->, line width=2.2pt, omDef] (2.0,1.7) -- node[pos=0.88, right=3pt, align=left, font=\scriptsize] {response $\lambda\,\Delta T$} (2.0,3.3);
\draw[->, line width=1.4pt, black!70] (0,-0.1) -- node[right=3pt, align=left, font=\scriptsize] {stored: $\Delta N = \Delta F + \lambda\,\Delta T$} (0,-1.0);
\draw[black!40, dashed] (-3.4,2.4) -- (3.4,2.4) node[right, font=\scriptsize, text=black!60] {top of the atmosphere};
\end{tikzpicture}
\par\smallskip{\small The energy budget of a forced climate (schematic). The forcing adds energy at
the top of the atmosphere, the response of the warmer planet removes some ($\lambda\,\Delta T$ is
negative), and the difference is stored, mostly in the ocean. Shares of the heat gained from 1971
to 2018.}
\end{omfigure}
% ledger: ar6ch7:inventory

\begin{method}[A climate sensitivity from the historical energy budget]\label{met:b2:anthropogenic-climate-change:energy-budget}
\begin{enumerate}
  \item Choose two periods, a reference close to equilibrium and a recent one; take the change in
        forcing $\Delta F$, in imbalance $\Delta N$ and in global temperature $\Delta T$.
  \item Solve the balance for the feedback parameter: $\lambda = (\Delta N - \Delta F)/\Delta T$.
  \item The sensitivity to a doubling is then $-F_{2\times}/\lambda$; the warming at the time of the
        forcing, scaled to $F_{2\times}$, is $F_{2\times}\Delta T/\Delta F$.
\end{enumerate}
\end{method}

From 1850--1900 to 2006--2019 the assessment takes $\Delta N = \qty{0.59}{W.m^{-2}}$, $\Delta T =
\qty{1.03}{\celsius}$ and $\Delta F = \qty{2.20}{W.m^{-2}}$. Then $\lambda = (0.59 - 2.20)/1.03 =
\qty{-1.56}{W.m^{-2}.\celsius^{-1}}$, and the sensitivity is $3.93/1.56 = \qty{2.5}{\celsius}$,
the value the assessment finds, with a very likely range of 1.6 to \qty{4.8}{\celsius}. But this is
an \emph{effective} sensitivity, for the feedbacks of the warming so far. The warming has been
strongest where feedbacks are most negative, in the western tropical Pacific, and weaker over the
eastern Pacific and the Southern Ocean; as that pattern evolves the feedback parameter is expected
to become less negative (\emph{high confidence}), so the long-term sensitivity is higher than the
historical estimate (\cref{exo:b2:anthropogenic-climate-change:7}).
% ledger: ar6ch7:energy-budget, ar6ch7:pattern, ipcc:erf-2xco2
% values: lambda_energy -1.563, ecs_energy 2.514 (anthropogenic.py numbers)

\section{Feedbacks and sensitivity}\label{sec:b2:anthropogenic-climate-change:feedbacks}

\begin{definition}[Planck response; water-vapour, lapse-rate and cloud feedbacks]\label{def:b2:anthropogenic-climate-change:feedbacks}
The feedback parameter is split into the contributions of the quantities that change with the
surface temperature. The \emph{Planck response}\index{Planck response} is the extra infrared
emitted to space when the surface and the whole troposphere warm by the same amount. The
\emph{water-vapour feedback}\index{water-vapour feedback} is the change in flux due to the increase
of atmospheric water vapour with temperature. The \emph{lapse-rate feedback}\index{lapse-rate
feedback} is the change in flux due to a warming that is not uniform with height. The \emph{cloud
feedback}\index{cloud feedback} is the change in flux due to the change in the amount, height and
brightness of clouds.
\end{definition}

\begin{proposition}[The Planck response]\label{prop:b2:anthropogenic-climate-change:planck}
If the Earth emits to space a fixed fraction $1 - g$ of the infrared $\sigma T^4$ of its surface at
temperature $T$, its \omterm{def:b2:anthropogenic-climate-change:feedbacks}{Planck response} is
\[
  \lambda_{\text P} = -(1 - g)\,4\sigma T^3 .
\]
With $g = 0.4$ and $T = \qty{288}{K}$, $\lambda_{\text P} = \qty{-3.25}{W.m^{-2}.K^{-1}}$.
\end{proposition}
\begin{proof}
The outgoing flux is $(1 - g)\sigma T^4$; its derivative with respect to $T$ is $(1 - g)\,4\sigma
T^3$, and an extra loss is a negative contribution to the net downward flux. Numerically $0.6
\times 4 \times \num{5.670e-8} \times 288^3 = 3.25$.
\end{proof}
% ledger: ar6ch7:planck, const:sigma, ipcc:budget-surface
% values: planck -3.2509 (anthropogenic.py numbers)

The fraction $g$ is the normalised greenhouse effect: the atmosphere keeps \qty{159}{W.m^{-2}} of
the \qty{398}{W.m^{-2}} that the surface emits, $g \approx 0.4$. The assessment, with detailed
radiation codes, finds $-3.22$ [$-3.4$ to $-3.0$] (\emph{high confidence}). Water vapour follows
the Clausius--Clapeyron law at the global scale, and more vapour absorbs more infrared: the
\omterm{def:b2:anthropogenic-climate-change:feedbacks}{water-vapour feedback} is strongly positive, \qty{1.85}{W.m^{-2}.\celsius^{-1}} from satellites.
In the tropics the upper troposphere warms more than the surface, which then radiates more: the
\omterm{def:b2:anthropogenic-climate-change:feedbacks}{lapse-rate feedback} is negative there, positive near the poles, and about $-0.50$ globally. The two
are assessed together, \num{1.30} [1.1 to 1.5], because what one model puts in one it tends to
take from the other. Melting snow and ice darken the surface, +0.35; clouds add \num{0.42}
[$-0.10$ to $0.94$], the largest uncertainty. A net negative \omterm{def:b2:anthropogenic-climate-change:feedbacks}{cloud feedback} is \emph{very unlikely}.
% ledger: ar6ch7:planck, ar6ch7:wv-lr, ar6ch7:table710, ipcc:cloud-negative, ipcc:budget-surface

\begin{omfigure}
\centering
\begin{tikzpicture}
\begin{axis}[xbar, width=0.78\linewidth, height=5.4cm, xmin=-3.7, xmax=1.8, ymin=0.4, ymax=6.6,
  ytick={1,...,6}, yticklabels={net, biogeophysical and chemical, clouds, surface albedo,
  water vapour + lapse rate, Planck}, y tick label style={font=\scriptsize},
  tick label style={font=\scriptsize}, xlabel={feedback parameter (\unit{W.m^{-2}.\celsius^{-1}})},
  label style={font=\small}, bar width=7pt, axis x line=bottom, axis y line=left,
  axis line style={-}, xmajorgrids, grid style={gray!22, line width=0.3pt}, extra x ticks={0},
  extra x tick style={grid style={black!60, line width=0.5pt}}]
% ledger: ar6ch7:table710 (central values and very likely ranges)
\addplot[fill=omDef!55, draw=omDef!80!black, error bars/.cd, x dir=both, x explicit,
  error bar style={black!70}] coordinates {
  (-3.22,6) += (0.22,0) -= (0.18,0)
  (-0.01,2) += (0.26,0) -= (0.26,0)
  (-1.16,1) += (0.65,0) -= (0.65,0)};
\addplot[fill=omThm!55, draw=omThm!80!black, error bars/.cd, x dir=both, x explicit,
  error bar style={black!70}] coordinates {
  (1.30,5) += (0.20,0) -= (0.20,0)
  (0.35,4) += (0.25,0) -= (0.25,0)
  (0.42,3) += (0.52,0) -= (0.52,0)};
\end{axis}
\end{tikzpicture}
\par\smallskip{\small The climate feedbacks assessed by the IPCC (Sixth Assessment Report, Working
Group~I, Table 7.10), central values and very likely ranges. Positive feedbacks amplify a warming
(red), negative ones damp it (blue); the net is the sum.}
\end{omfigure}

\begin{definition}[Equilibrium climate sensitivity]\label{def:b2:anthropogenic-climate-change:ecs}
The \emph{equilibrium climate sensitivity}\index{equilibrium climate sensitivity} (ECS) is the
global surface warming, at the new equilibrium, caused by a doubling of the carbon dioxide of the
air from its pre-industrial level, the ice sheets being held fixed.
\end{definition}

\begin{proposition}[Sensitivity and feedback]\label{prop:b2:anthropogenic-climate-change:ecs}
If the feedback parameter is constant, $\text{ECS} = -F_{2\times}/\lambda$. With $F_{2\times} =
\qty{3.93}{W.m^{-2}}$ and the net $\lambda = \qty{-1.16}{W.m^{-2}.\celsius^{-1}}$ of the feedbacks
above, $\text{ECS} = \qty{3.4}{\celsius}$; the \omterm{def:b2:anthropogenic-climate-change:feedbacks}{Planck response} alone would give $3.93/3.22 =
\qty{1.2}{\celsius}$.
\end{proposition}
\begin{proof}
Apply \cref{prop:b2:anthropogenic-climate-change:balance} at equilibrium with $\Delta F =
F_{2\times}$.
\end{proof}
% ledger: ar6ch7:table710, ipcc:erf-2xco2
% values: ecs_process 3.388 (anthropogenic.py numbers)

The feedbacks are one line of evidence; the warming since 1850, the energy budget and the climates
of the past are others. Combined, they give the assessed ECS of \qty{3}{\celsius}, \emph{likely}
2.5 to \qty{4}{\celsius} (\emph{high confidence}), \emph{very likely} 2 to \qty{5}{\celsius}.
Because the ECS varies as $1/\lambda$, an uncertainty on $\lambda$ that is symmetric gives a range
of ECS stretched towards high values (\cref{exo:b2:anthropogenic-climate-change:12}).
% ledger: ipcc:ecs

\begin{definition}[Transient climate response]\label{def:b2:anthropogenic-climate-change:tcr}
The \emph{transient climate response}\index{transient climate response} (TCR) is the global surface
warming at the moment the carbon dioxide of the air has doubled, when it rises by \qty{1}{\percent}
a year from its pre-industrial level, about seventy years after the start.
\end{definition}

The TCR is smaller than the ECS because the ocean is still taking up heat when the doubling is
reached. Its best estimate is \qty{1.8}{\celsius}, \emph{likely} 1.4 to \qty{2.2}{\celsius}
(\emph{high confidence}); the historical energy budget alone gives $3.93 \times 1.03/2.20 =
\qty{1.8}{\celsius}$ as well.
% ledger: ar6ch7:tcr, ar6ch7:energy-budget
% values: tcr_energy 1.840 (anthropogenic.py numbers)

\section{Detection and attribution}\label{sec:b2:anthropogenic-climate-change:attribution}

\begin{definition}[Detection; attribution; climate fingerprint]\label{def:b2:anthropogenic-climate-change:da}
\emph{Detection}\index{detection (of climate change)} of a climate change is showing that the
climate has changed in a defined statistical sense, beyond what its internal variability could
produce, without giving the cause. \emph{Attribution}\index{attribution} is evaluating the
contributions of several causes to a detected change, with an assigned likelihood. A \emph{climate
fingerprint}\index{climate fingerprint} is the pattern in space and time of the response to one
driver, as simulated by climate models, by which that driver is recognised in the observations.
\end{definition}

Different drivers leave different fingerprints. More greenhouse gas warms the troposphere and cools
the stratosphere, because a thicker infrared blanket makes the upper layers radiate more to space;
a brighter Sun would warm both. The observations show the first pattern: the troposphere has warmed
since at least the 1950s, and it is \emph{virtually certain} that the stratosphere has cooled. It is
\emph{very likely} that the human greenhouse gases were the main driver of the tropospheric warming
since 1979, and \emph{extremely likely} that the depletion of the ozone layer was the main driver of
the cooling of the lower stratosphere from 1979 to the mid-1990s.
% ledger: ar6ch3:upper-air

\begin{method}[Optimal fingerprinting]\label{met:b2:anthropogenic-climate-change:fingerprinting}
\begin{enumerate}
  \item Simulate the response to each driver separately: the fingerprints $x_1, \dots, x_k$, as
        vectors of values in space and time.
  \item Write the observations as $y = \sum_i \beta_i x_i + \varepsilon$, with $\varepsilon$ the
        internal variability, whose statistics are taken from long unforced simulations.
  \item Estimate the scaling factors $\beta_i$ by least squares, weighted by the inverse of the
        variability so that the noisiest components count least, with their confidence intervals.
  \item A $\beta_i$ significantly greater than zero means that the response to driver $i$ is
        detected; an interval that contains 1 and excludes 0 means that the simulated response is
        consistent with the observations, and the change can be attributed in part to that driver.
\end{enumerate}
\end{method}
% ledger: ar6ch3:fingerprint

Applied to the surface temperature, the method gives the \omterm{def:b2:anthropogenic-climate-change:da}{attribution} recalled at the start: of the
\qty{1.06}{\celsius} observed between 1850--1900 and 2010--2019, the human influence \emph{likely}
caused 0.8 to \qty{1.3}{\celsius}, best estimate \qty{1.07}{\celsius}; the natural drivers and the
internal variability contribute no more than a few tenths of a degree either way
(\cref{exo:b2:anthropogenic-climate-change:9,exo:b2:anthropogenic-climate-change:10}).
% ledger: ipcc:attribution, ipcc:observed-2010-2019

\section{Scenarios and cumulative emissions}\label{sec:b2:anthropogenic-climate-change:scenarios}

\begin{definition}[Shared Socioeconomic Pathway]\label{def:b2:anthropogenic-climate-change:ssp}
A \emph{Shared Socioeconomic Pathway}\index{Shared Socioeconomic Pathway} (SSP) is one of five
narratives of how population, economies and technology might develop over the century. The
emission scenarios of the assessment are labelled SSP$x$-$y$: $x$ names the pathway, and $y$ the
approximate radiative forcing, in \unit{W.m^{-2}}, that the scenario reaches in 2100.
\end{definition}
% ledger: ar6spm:ssp

The five main scenarios run from SSP1-1.9 to SSP5-8.5. The physical science assessment takes them as
given: it neither weighs their likelihood nor assesses the narratives. For each, the climate models
give a projected warming for 2081--2100 above 1850--1900, from \qty{1.4}{\celsius} (very likely 1.0
to 1.8) under SSP1-1.9 to \qty{4.4}{\celsius} (3.3 to 5.7) under SSP5-8.5. The projections share a
simple law.
% ledger: ar6spm:ssp, ipcc:proj-temp

\begin{definition}[Transient climate response to cumulative CO$_2$ emissions]\label{def:b2:anthropogenic-climate-change:tcre}
The \emph{transient climate response to cumulative \ce{CO2} emissions}\index{transient climate
response to cumulative CO$_2$ emissions} (TCRE) is the global surface warming per unit of carbon
dioxide emitted in total since the pre-industrial period, in \unit{\celsius} per
\qty{1000}{Gt} of \ce{CO2}.
\end{definition}

\begin{proposition}[Warming follows cumulative emissions]\label{prop:b2:anthropogenic-climate-change:tcre}
To a good approximation the warming caused by carbon dioxide is proportional to the total emitted,
whatever the path of the emissions: $\Delta T \approx \text{TCRE} \times E$, with $E$ the cumulative
emissions.
\end{proposition}
\admitted
% ledger: ar6spm:tcre, ar6ch5:tcre-linear

The law is not obvious. The forcing grows only as the logarithm of the concentration, so each
added tonne warms less; but the ocean takes up heat and carbon less and less efficiently as
emissions accumulate, so more of each tonne stays in the air and its warming is less delayed. The
two effects nearly cancel, which all the studies agree is the origin of the law. The relationship is
near-linear (\emph{high confidence}), and each \qty{1000}{Gt} of carbon dioxide \emph{likely} causes
0.27 to \qty{0.63}{\celsius} of warming, best estimate \qty{0.45}{\celsius}, or \qty{1.65}{\celsius}
per \qty{1000}{PgC} of carbon (\cref{exo:b2:anthropogenic-climate-change:4}). From 1850 to 2019,
$\num{2390} \pm \qty{240}{Gt}$ of carbon dioxide were emitted.
% ledger: ar6ch5:tcre-linear, ar6spm:tcre, ar6spm:budget

\begin{definition}[Remaining carbon budget]\label{def:b2:anthropogenic-climate-change:budget}
The \emph{remaining carbon budget}\index{remaining carbon budget} for a temperature limit is the
cumulative carbon dioxide that can still be emitted, from a stated date, with global warming staying
below that limit at a stated probability, the warming caused by the other human drivers being taken
into account.
\end{definition}

\begin{method}[The remaining carbon budget from the TCRE]\label{met:b2:anthropogenic-climate-change:budget}
\begin{enumerate}
  \item Subtract the human warming already reached, \qty{1.07}{\celsius} for 2010--2019, from the
        limit.
  \item Subtract the warming the other drivers will add by the time carbon dioxide emissions
        reach net zero.
  \item Divide by the TCRE: this is the budget from the middle of 2010--2019.
  \item Subtract the emissions since then, and a margin for the slow Earth-system feedbacks the
        models leave out (permafrost thaw is the largest).
\end{enumerate}
\end{method}

The assessment combines five components in this way: the TCRE; the historical warming; an
additional warming by the other drivers of about 0.1 to \qty{0.2}{\celsius}; the warming that may
continue after emissions stop, assessed at zero on average (\emph{likely} within $\pm
\qty{0.19}{\celsius}$); and the Earth-system feedbacks, about \qty{26}{Gt} of carbon dioxide per
degree. For \qty{1.5}{\celsius}, \qty{0.43}{\celsius} above 2010--2019, the remaining budget from the
start of 2020 is \qty{500}{Gt} of carbon dioxide at a \qty{50}{\percent} probability, and
\qty{1350}{Gt} for \qty{2}{\celsius}. The choices made about the other drivers move these values by
\qty{220}{Gt} or more; the weekend problem rebuilds the first one.
% ledger: ar6ch5:budget-parts, ar6spm:budget

\begin{dated}{2025}
% ledger: data:co2-mlo, data:gistemp (values read from the datasets as fetched 2026-10-22; offset computed)
The annual mean carbon dioxide at Mauna Loa reached \qty{427.35}{ppm} in 2025, a forcing of
$3.93\log_2(427.35/278) = \qty{2.44}{W.m^{-2}}$ by the logarithmic law. In the NASA temperature
record 2024 is the warmest year, \qty{1.50}{\celsius} above the mean of 1880--1900, the earliest
years of that record; 2025 stood at \qty{1.40}{\celsius}.
\end{dated}

\begin{omfigure}
\centering
\begin{tikzpicture}
\begin{axis}[omrecord, width=0.86\linewidth, height=5.8cm, xmin=0, xmax=4000, ymin=0, ymax=2.6,
  xlabel={cumulative \ce{CO2} emissions since 1850 (\unit{Gt})},
  ylabel={warming since 1850--1900 (\unit{\celsius})}, xtick={0,500,...,4000},
  /pgf/number format/1000 sep={\,}, xtick pos=bottom, ytick pos=left, clip=false,
  tick label style={font=\scriptsize}, label style={font=\small}]
% computed in figdata/bachelor-2/anthropogenic.py part tcre; ledger: ar6spm:tcre, ar6spm:budget
\addplot[name path=hi, draw=none, forget plot] table[x=E_GtCO2, y=dT_hi] {figdata/out/bachelor-2-anthropogenic-tcre.dat};
\addplot[name path=lo, draw=none, forget plot] table[x=E_GtCO2, y=dT_lo] {figdata/out/bachelor-2-anthropogenic-tcre.dat};
\addplot[omThm!18, forget plot] fill between[of=hi and lo];
\addplot[omThm, line width=1.3pt] table[x=E_GtCO2, y=dT_mid] {figdata/out/bachelor-2-anthropogenic-tcre.dat};
\draw[densely dashed, black!60] (axis cs:0,1.5) -- (axis cs:4000,1.5);
\draw[densely dashed, black!60] (axis cs:0,2.0) -- (axis cs:4000,2.0);
\node[anchor=north west, font=\scriptsize, text=black!70] at (axis cs:0,1.5) {\qty{1.5}{\celsius}};
\node[anchor=north west, font=\scriptsize, text=black!70] at (axis cs:0,2.0) {\qty{2}{\celsius}};
\fill[black] (axis cs:2390,1.07) circle (2.2pt);
\draw[black, line width=0.8pt] (axis cs:2150,1.07) -- (axis cs:2630,1.07);
\draw[black, line width=0.8pt] (axis cs:2390,0.8) -- (axis cs:2390,1.3);
\node[anchor=north west, font=\scriptsize, align=left] at (axis cs:2440,1.02) {1850--2019:\\\qty{2390}{Gt}, \qty{1.07}{\celsius}};
\draw[<->, omDef, thick] (axis cs:2390,1.62) -- (axis cs:2890,1.62);
\node[anchor=south, font=\scriptsize, text=omDef, align=center] at (axis cs:2640,1.64) {AR6: \qty{500}{Gt}};
\node[anchor=north west, font=\scriptsize, text=omThm, align=left] at (axis cs:150,2.5) {likely range of the TCRE,\\0.27 to \qty{0.63}{\celsius} per \qty{1000}{Gt}};
\end{axis}
\end{tikzpicture}
\par\smallskip{\small Warming against cumulative carbon dioxide emissions for the assessed TCRE,
\qty{0.45}{\celsius} per \qty{1000}{Gt} (line) and its likely range (shaded), with the human warming
and the emissions to date (dot, likely ranges as bars). The arrow marks the remaining budget for
\qty{1.5}{\celsius} at \qty{50}{\percent} from 2020: smaller than a reading of the line alone,
because the other drivers also warm.}
\end{omfigure}


\section{Sea level}\label{sec:b2:anthropogenic-climate-change:sea-level}

\begin{definition}[Sea-level budget; thermosteric sea-level rise]\label{def:b2:anthropogenic-climate-change:sea-level}
The \emph{sea-level budget}\index{sea-level budget} is the comparison of the observed change of
global mean sea level with the sum of its estimated contributions: the expansion of the warming
ocean, the melting of glaciers and of the Greenland and Antarctic ice sheets, and the change of the
water stored on land. \emph{Thermosteric sea-level rise}\index{thermosteric sea-level rise} is the
part due to the thermal expansion of sea water, at constant mass.
\end{definition}

\begin{proposition}[Thermal expansion of a layer]\label{prop:b2:anthropogenic-climate-change:thermosteric}
A layer of ocean of thickness $H$, with thermal expansion coefficient $\alpha$, that warms uniformly
by $\Delta T$ at constant mass and area rises by $\Delta h = \alpha\,\Delta T\,H$.
\end{proposition}
\begin{proof}
The mass $\rho H A$ of a column of area $A$ is constant, so $\Delta H/H = -\Delta\rho/\rho =
\alpha\,\Delta T$ to first order, with $\alpha = -(1/\rho)\,\partial\rho/\partial T$.
\end{proof}
% ledger: teos10:alpha-beta

In the real ocean $\alpha$ grows with temperature and pressure, and the warming is concentrated near
the surface. The assessment summarises the result as an expansion efficiency of \qty{0.113}{m} per
\qty{e24}{J} of heat added to the ocean, or \qty{0.617}{m} per degree of warming of the whole ocean
(\emph{medium confidence}). Over its mean depth of \qty{3682}{m}, that is an effective $\alpha$ of
$0.617/3682 = \qty{1.68e-4}{\celsius^{-1}}$ (\cref{exo:b2:anthropogenic-climate-change:8}).
% ledger: ar6ch9:expansion, charette:ocean-volume

The budget, drawn in the figure below, closes within its uncertainties. Over 1971--2018 the contributions add up to
\qty{2.00}{mm.yr^{-1}}, against an observed rise of \qty{2.33}{mm.yr^{-1}}; over 2006--2018 they
add up to 3.61 against 3.69. Thermal expansion was the largest single term, half of the sum over 1971--2018, but the ice is
catching up: over 2006--2018 glaciers and ice sheets together gave \qty{1.62}{mm.yr^{-1}}, more than
the expansion. The ocean also keeps warming at depth long after the surface settles, and the ice
sheets respond over centuries; sea level is committed to rise for centuries to millennia and will
remain high for thousands of years (\emph{high confidence}). Over the next two thousand years it is
projected to rise by about 2 to \qty{3}{m} if the warming stays at \qty{1.5}{\celsius}, 2 to
\qty{6}{m} at \qty{2}{\celsius} and 19 to \qty{22}{m} at \qty{5}{\celsius} (\emph{low confidence}).
% ledger: ar6ch9:table95, ar6spm:sea-commit

\begin{omfigure}
\centering
\begin{tikzpicture}
\begin{axis}[ybar stacked, width=0.72\linewidth, height=6.2cm, ymin=0, ymax=4.3, bar width=26pt,
  xtick=data, symbolic x coords={1971--2018, 1993--2018, 2006--2018}, enlarge x limits=0.25,
  ylabel={rate of rise (\unit{mm.yr^{-1}})}, tick label style={font=\scriptsize},
  label style={font=\small}, ytick pos=left, xtick pos=bottom,
  legend style={at={(1.03,0.5)}, anchor=west, font=\scriptsize, draw=black!30},
  legend cell align=left, reverse legend]
% ledger: ar6ch9:table95
\addplot[fill=omDef!70, draw=omDef!90!black] coordinates {(1971--2018,1.01) (1993--2018,1.31) (2006--2018,1.39)};
\addlegendentry{thermal expansion}
\addplot[fill=omDef!30, draw=omDef!70!black] coordinates {(1971--2018,0.44) (1993--2018,0.55) (2006--2018,0.62)};
\addlegendentry{glaciers}
\addplot[fill=omThm!55, draw=omThm!80!black] coordinates {(1971--2018,0.25) (1993--2018,0.43) (2006--2018,0.63)};
\addlegendentry{Greenland}
\addplot[fill=omThm!25, draw=omThm!70!black] coordinates {(1971--2018,0.14) (1993--2018,0.25) (2006--2018,0.37)};
\addlegendentry{Antarctica}
\addplot[fill=black!20, draw=black!50] coordinates {(1971--2018,0.15) (1993--2018,0.31) (2006--2018,0.60)};
\addlegendentry{land water}
\node[fill=black, rotate=45, minimum size=5pt, inner sep=0pt] at (axis cs:1971--2018,2.33) {};
\node[fill=black, rotate=45, minimum size=5pt, inner sep=0pt] at (axis cs:1993--2018,3.25) {};
\node[fill=black, rotate=45, minimum size=5pt, inner sep=0pt] at (axis cs:2006--2018,3.69) {};
\end{axis}
\end{tikzpicture}
\par\smallskip{\small The \omterm{def:b2:anthropogenic-climate-change:sea-level}{sea-level budget}: observed contributions to the rise of global mean sea
level, as rates over three periods ending in 2018, stacked, and the observed rate (diamonds). The
gap between the stack and the diamond is the part the budget does not close. Values: IPCC Sixth
Assessment Report, Working Group~I, Table 9.5.}
\end{omfigure}

\section*{Exercises}\label{sec:b2:anthropogenic-climate-change:exercises}

Data for the exercises, unless stated otherwise: forcing of doubled carbon dioxide $F_{2\times} =
\qty{3.93}{W.m^{-2}}$ from $C_0 = \qty{278}{ppm}$; Stefan--Boltzmann constant $\sigma =
\qty{5.670e-8}{W.m^{-2}.K^{-4}}$; radius of the Earth \qty{6371}{km}; one year
\qty{3.156e7}{s}; molar masses of carbon and carbon dioxide 12.011 and \qty{44.009}{g.mol^{-1}};
the feedbacks and the \omterm{def:b2:anthropogenic-climate-change:sea-level}{sea-level budget} of the chapter.
% ledger: ipcc:erf-2xco2, const:sigma, earth:Rmean, ciaaw:abridged-2024, ar6ch7:table710, ar6ch9:table95

\begin{exercise}[$\star$]\label{exo:b2:anthropogenic-climate-change:1}
The forcing of doubled carbon dioxide is \qty{3.75}{W.m^{-2}} once the stratosphere has adjusted;
the troposphere then adjusts by $-0.60$ (temperature), +0.22 (water vapour), +0.45 (clouds) and
+0.11 (land surface) \unit{W.m^{-2}}. (a) Find the ERF. (b) By what percentage does it exceed the
previous assessment's \qty{3.71}{W.m^{-2}}? (c) Which adjustment is negative, and why does a warmer
troposphere reduce the forcing?
\end{exercise}
% ledger: ar6ch7:2xco2

\begin{exercise}[$\star$]\label{exo:b2:anthropogenic-climate-change:2}
The climate system gained \qty{434.9}{ZJ} from 1971 to 2018 and \qty{152.4}{ZJ} from 2006 to 2018,
counted between the middles of the first and last years. (a) Find the Earth's surface area.
(b) Convert each gain into a mean imbalance in \unit{W.m^{-2}}. (c) The ocean took
\qty{91}{\percent}: how many zettajoules over 1971--2018?
\end{exercise}
% ledger: ar6ch7:inventory, earth:Rmean

\begin{exercise}[$\star$]\label{exo:b2:anthropogenic-climate-change:3}
(a) Add up the central values of the feedbacks in the figure, and compare with the net. (b) Find
the ECS with the net value, and the warming with the \omterm{def:b2:anthropogenic-climate-change:feedbacks}{Planck response} alone. (c) By what factor do
the other feedbacks amplify the Planck warming?
\end{exercise}

\begin{exercise}[$\star$]\label{exo:b2:anthropogenic-climate-change:4}
(a) Convert a TCRE of \qty{0.45}{\celsius} per \qty{1000}{Gt} of carbon dioxide into degrees per
\qty{1000}{PgC} of carbon, and do the same for its likely range, 0.27 to 0.63. (b) Convert the
\qty{2390}{Gt} of carbon dioxide emitted from 1850 to 2019 into petagrams of carbon. (c) What
warming does the TCRE give for these emissions?
\end{exercise}
% ledger: ar6spm:tcre, ar6spm:budget

\begin{exercise}[$\star\star$]\label{exo:b2:anthropogenic-climate-change:5}
(a) The surface emits \qty{398}{W.m^{-2}} and the greenhouse effect keeps \qty{159}{W.m^{-2}} of
it. Find $g$ and the surface temperature of a black body emitting \qty{398}{W.m^{-2}}. (b) Find the
\omterm{def:b2:anthropogenic-climate-change:feedbacks}{Planck response} at \qty{288}{K}, and at \qty{298}{K} with the same $g$. (c) Would the \omterm{def:b2:anthropogenic-climate-change:feedbacks}{Planck
response} alone make a warmer Earth more or less sensitive?
\end{exercise}
% ledger: ipcc:budget-surface, ar6ch7:planck

\begin{exercise}[$\star\star$]\label{exo:b2:anthropogenic-climate-change:6}
Satellites give a \omterm{def:b2:anthropogenic-climate-change:feedbacks}{water-vapour feedback} of $1.85 \pm 0.32$ and the studies a \omterm{def:b2:anthropogenic-climate-change:feedbacks}{lapse-rate feedback} of
$-0.50 \pm 0.20$ \unit{W.m^{-2}.\celsius^{-1}} (one standard deviation). (a) Find their sum.
(b) If the two were independent, what would be the standard deviation of the sum? (c) The models
give the sum a standard deviation of only \num{0.10}. What does that say about the two errors?
\end{exercise}
% ledger: ar6ch7:wv-lr

\begin{exercise}[$\star\star$]\label{exo:b2:anthropogenic-climate-change:7}
From 1850--1900 to 2006--2019, $\Delta N = \qty{0.59}{W.m^{-2}}$, $\Delta T = \qty{1.03}{\celsius}$
and $\Delta F = \qty{2.20}{W.m^{-2}}$. (a) Find $\lambda$ and the effective sensitivity. (b) The
feedback parameter is expected to become less negative by $\lambda' = \qty{0.5}{W.m^{-2}.\celsius^{-1}}$
as the pattern of warming evolves. Find the sensitivity $-F_{2\times}/(\lambda + \lambda')$.
(c) Redo (b) with $\lambda' = 0.1$.
\end{exercise}
% ledger: ar6ch7:energy-budget, ar6ch7:pattern, ar6ch7:pattern-ecs

\begin{exercise}[$\star\star$]\label{exo:b2:anthropogenic-climate-change:8}
(a) From the \qty{0.617}{m} of rise per degree of warming of the whole ocean, mean depth
\qty{3682}{m}, find the effective expansion coefficient. (b) The ocean gained \qty{396}{ZJ} over
1971--2018. With an expansion efficiency of \qty{0.113}{m} per \qty{e24}{J}, find the thermosteric
rise, and compare with the \qty{47.5}{mm} of the budget. (c) A layer of \qty{700}{m} warms by
\qty{0.2}{\celsius} with $\alpha = \qty{2e-4}{\celsius^{-1}}$: how much does it raise sea level?
\end{exercise}
% ledger: ar6ch9:expansion, charette:ocean-volume, ar6ch7:inventory, ar6ch9:table95

\begin{exercise}[$\star\star$]\label{exo:b2:anthropogenic-climate-change:9}
A model fingerprint, in four decades, is $x = (0.10, 0.20, 0.30, 0.40)$ \unit{\celsius}; the
observations are $y = (0.12, 0.18, 0.33, 0.41)$ \unit{\celsius}. (a) Find the least-squares scaling
factor $\beta$ of $y = \beta x + \varepsilon$. (b) Its standard error, estimated from unforced
simulations, is \num{0.08}. Is the response detected? Is it consistent with the model? (c) What
would a $\beta$ of 0.5 with the same error mean?
\end{exercise}

\begin{exercise}[$\star\star\star$]\label{exo:b2:anthropogenic-climate-change:10}
Two fingerprints, greenhouse gases $g = (0.2, 0.4, 0.6, 0.8)$ and aerosols $a = (-0.1, -0.2,
-0.2, -0.3)$, both in \unit{\celsius} over four decades, are fitted to the observations $y = (0.1,
0.2, 0.4, 0.5)$. (a) Write the normal equations of the least-squares fit $y = \beta_g g + \beta_a a
+ \varepsilon$ and solve them. (b) Find the warming attributed to each driver in the last decade.
(c) Why do the uncertainties on $\beta_g$ and $\beta_a$ grow when the two fingerprints look alike?
\end{exercise}

\begin{exercise}[$\star\star\star$]\label{exo:b2:anthropogenic-climate-change:11}
Use the method of the \omterm{def:b2:anthropogenic-climate-change:budget}{remaining carbon budget} with the TCRE of \qty{0.45}{\celsius} per
\qty{1000}{Gt}, the \qty{210}{Gt} emitted from 2015 to 2019 and \qty{26}{Gt} per degree for the
Earth-system feedbacks. (a) Find the budget for \qty{2}{\celsius}, \qty{0.93}{\celsius} above
2010--2019, if the other drivers add \qty{0.1}{\celsius}. (b) What additional warming by the other
drivers makes it equal to the assessed \qty{1350}{Gt}? Is that within the assessed range? (c) Show
that each \qty{0.1}{\celsius} of non-CO$_2$ warming changes any budget by the same amount.
\end{exercise}

\begin{exercise}[$\star\star\star$]\label{exo:b2:anthropogenic-climate-change:12}
The net feedback has a very likely range of $-1.81$ to \qty{-0.51}{W.m^{-2}.\celsius^{-1}}, around
$-1.16$. (a) Find the ECS at the two ends and at the centre. (b) Show that a range of $\lambda$
symmetric about its centre gives an ECS range that is not, and say which side is longer. (c) Why
do the assessments find it harder to rule out a high ECS than a low one?
\end{exercise}
% ledger: ar6ch7:table710

\begin{problem}[{Weekend problem --- Reading the budget}]\label{pb:b2:anthropogenic-climate-change:1}
This problem follows the energy and the carbon of the assessment from the forcing to the remaining
budget for \qty{1.5}{\celsius}. Use the data of the exercises, and: ERF 1750--2019 (\unit{W.m^{-2}})
carbon dioxide 2.16, methane 0.54, nitrous oxide 0.21, halogenated gases 0.41, ozone 0.47,
stratospheric water vapour 0.05, aerosol--radiation $-0.22$, aerosol--cloud $-0.84$, land use
$-0.20$, black carbon on snow 0.08, contrails 0.06; carbon dioxide \qty{410}{ppm} in 2019 and
\qty{427.35}{ppm} in 2025; heating \qty{0.79}{W.m^{-2}} over 2006--2018; the energy-budget changes
$\Delta N = 0.59$, $\Delta F = \qty{2.20}{W.m^{-2}}$, $\Delta T = \qty{1.03}{\celsius}$; TCRE
\qty{0.45}{\celsius} per \qty{1000}{Gt} of carbon dioxide (likely 0.27 to 0.63); \qty{2390}{Gt}
emitted 1850--2019; human warming \qty{1.07}{\celsius} for 2010--2019; \qty{210}{Gt} emitted from
2015 to 2019; additional non-CO$_2$ warming \qty{0.1}{\celsius}, the lower end of the assessed 0.1
to \qty{0.2}{\celsius}; Earth-system feedbacks \qty{26}{Gt} per degree. Round to three significant
figures.
% ledger: ipcc:erf-agents, ipcc:erf-2xco2, data:co2-mlo, ar6ch7:inventory, ar6ch7:energy-budget, ar6spm:tcre, ar6spm:budget, ar6ch5:budget-parts

\noindent\textbf{Part I --- The forcing.}
\begin{enumerate}
  \item Add up the forcing of the four well-mixed greenhouse gases.
  \item Add up the forcing of the aerosols.
  \item Add up all the human terms. Compare with the assessed \qty{2.72}{W.m^{-2}}.
  \item Check the forcing of carbon dioxide in 2019 with the logarithmic law.
  \item Find the forcing of carbon dioxide in 2025 by the same law.
  \item Which term has the widest uncertainty, and why does it dominate that of the total?
\end{enumerate}

\noindent\textbf{Part II --- The energy budget.}
\begin{enumerate}[resume]
  \item How much energy did the climate system gain over 2006--2018, in zettajoules?
  \item How much of it went into the ocean?
  \item Find the feedback parameter from the energy-budget changes.
  \item Find the effective \omterm{def:b2:anthropogenic-climate-change:ecs}{equilibrium climate sensitivity}.
  \item Find the \omterm{def:b2:anthropogenic-climate-change:tcr}{transient climate response} from the same data.
  \item Find the ECS from the net feedback of the figure, and compare the two values.
  \item Give the physical reason the historical estimate is the lower one.
\end{enumerate}

\noindent\textbf{Part III --- Warming per tonne.}
\begin{enumerate}[resume]
  \item Convert the TCRE to degrees per \qty{1000}{PgC} of carbon.
  \item Convert its likely range.
  \item Convert the emissions of 1850--2019 into petagrams of carbon.
  \item What warming does the TCRE give for the emissions of 1850--2019?
  \item Compare with the human warming of \qty{1.07}{\celsius}. What does the comparison say of the
        net effect of the other human drivers?
\end{enumerate}

\noindent\textbf{Part IV --- The remaining budget for \qty{1.5}{\celsius}.}
\begin{enumerate}[resume]
  \item What warming is left above 2010--2019?
  \item Remove the non-CO$_2$ warming and divide by the TCRE: what budget remains from the middle
        of 2010--2019?
  \item Remove the emissions of 2015--2019.
  \item Remove the Earth-system feedbacks over the remaining warming.
  \item By how much does the budget change if the other drivers add \qty{0.2}{\celsius} instead?
  \item To sum up: what is the \omterm{def:b2:anthropogenic-climate-change:budget}{remaining carbon budget} for \qty{1.5}{\celsius}, from the start of
        2020, computed from the TCRE? Compare with the assessed \qty{500}{Gt}.
\end{enumerate}
\end{problem}
